\section{Descriptive Analysis Plan}

\subsection{Planned Descriptive Statistics and Their Application}

The primary analysis will consist of descriptive statistical summaries of the survey respondents and their reported characteristics, behaviours, preferences, and experiences. The statistical summary applied to each variable will be determined by its measurement type and distribution. The analysis will prioritise description of the observed sample, with descriptive estimates remaining the principal analytical outputs.

Categorical and binary variables will generally be summarised using frequencies and percentages. Numerical variables will be summarised using appropriate measures of central tendency and dispersion, including the mean and standard deviation and/or median and interquartile range where appropriate. Ordinal and count variables will be summarised using distributions and appropriate measures of central tendency and dispersion according to their measurement properties.

% TBD: The final statistical software and precise reporting conventions will be established before analysis.

\subsection{Primary Variables and Outcomes}

The principal variables of interest will describe four broad aspects of the surveyed Fimfiction community: \textbf{who respondents are, how they engage with the My Little Pony fandom, how they engage with Fimfiction.net, and what types of content and activities they prefer}.

Demographic variables will include age, country of residence, highest educational attainment, political orientation, religious affiliation, disability, neurodivergence, LGBTQ+ identity, and other relevant characteristics. Fandom variables will include fandom entry and history, age at entry, generation engagement, G6 interest, fandom identification, engagement with other platforms, convention participation, social relationships formed through the fandom, and participation in non-MLP activities with other fans.

Fimfiction-specific variables will include site and group engagement, reading behaviour, story-rating preferences, tag preferences, character preferences, romance and pairing preferences, and interest in different types of characters and generations. Additional variables will describe participation as a reader and/or writer and, for writers, the use of generative AI for creative and related activities.

% TBD: The final designation of primary versus secondary variables will be established following finalisation
% of the survey instrument and research questions.

\subsection{Specification of Descriptive Summaries}

Categorical and binary variables will be presented as counts and percentages for each response state. For select-all-that-apply questions, each response category will be treated as an individual binary variable, with the frequency and percentage of respondents selecting that category reported separately.

Ordinal variables will generally be summarised using frequencies and percentages for each response category and may additionally be summarised using the median and interquartile range where appropriate. Numerical variables, including age, age at fandom entry, and 1--10 preference ratings, will be summarised using appropriate measures of central tendency, dispersion, and distribution.

Count variables may include the number of Fimfiction groups joined, other fandoms engaged with, platforms used for MLP participation, conventions attended, and other forms of community participation. These variables will be summarised using frequency distributions and appropriate measures of central tendency and dispersion.

Free-text responses may receive descriptive categorisation and frequency summaries where responses are coded into analytical categories.

% TBD: The final variable-by-variable mapping of descriptive summaries will be completed once the survey
% instrument, response options, and coding scheme have been finalised.

\subsection{Planned Descriptive Estimates and Denominators}

For categorical and binary variables, estimates will primarily be reported as counts and percentages. The denominator will generally consist of respondents with a valid substantive response to the variable being described. For conditional questions, the denominator will generally consist of respondents eligible to receive and answer the question rather than the entire survey sample.

Numerical and ordinal variables will be summarised using appropriate measures of central tendency, dispersion, and distribution. Select-all-that-apply categories will be reported independently, and percentages may therefore sum to greater than 100\%.

% TBD: The precise denominator rules for each conditional question will be specified once the final survey
% branching structure has been established.

\subsection{Distributional Assessment and Visualisation}

The distributions of numerical, ordinal, and count variables will be examined before determining the most appropriate descriptive summaries. Assessment will consider the range and concentration of observations, skewness, floor and ceiling effects, outliers, and other distributional characteristics relevant to interpretation.

Visual methods may include histograms, box plots, bar charts, density plots, or other appropriate graphical representations. Visual assessment may also identify unusual patterns that warrant further data-quality review or exploratory analysis.

% TBD: The final set of distributional diagnostics and visualisation methods will be specified before analysis.

\subsection{Overall Sample Description}

The overall sample will be described using a structured summary of demographic characteristics, fandom history, Fimfiction engagement, content preferences, social participation, and creative activity. The description will report the number of eligible and analysed respondents and, where relevant, the number excluded during data cleaning.

The characteristics of the observed respondents will be distinguished from claims about the wider Fimfiction.net user population.

\subsection{Statistical Thresholds and Multiple Comparisons}

The primary descriptive analyses will not depend on inferential statistical hypothesis testing. Descriptive estimates will be interpreted on the basis of their magnitude, distribution, uncertainty where applicable, and substantive relevance rather than according to statistical significance.

No statistical significance threshold or correction for multiple comparisons is planned for the primary descriptive analyses. If inferential statistical tests are subsequently included as supplementary or exploratory analyses, the relevant statistical thresholds, tests, and multiple-comparison procedures will be specified before those analyses are conducted.

% TBD: Whether inferential statistical tests will be included as supplementary or exploratory analyses, and any
% associated significance thresholds or multiple-comparison procedures, will be determined before those analyses.
